\section{Limitations}

Our work has several limitations that should be acknowledged. First, the human evaluation reveals moderate inter-annotator agreement ($\kappa=0.162$), reflecting the inherently subjective nature of ELI5 quality assessment. While both evaluators agree on accuracy and completeness improvements, they diverge on simplicity, with Evaluator 1 rating multi-agent outputs as less simple ($-$10.8\%) while Evaluator 2 rates them as more simple (+3.7\%). This divergence has not been resolved through qualitative annotation, and the rubric's definition of "simplicity for a 5-year-old" may need refinement to ensure consistent interpretation. Future work should include a small annotation study to calibrate evaluators on what constitutes appropriate simplicity in this context.

Second, we do not conduct ablation studies to isolate the contribution of individual agents. While the multi-agent architecture consistently outperforms single-pass prompting, we cannot attribute gains to specific components (e.g., the breakdown agent's decomposition vs.\ the synthesis agent's adaptive strategy). Future work should systematically remove or modify individual agents to quantify their marginal contribution.

Third, our evaluation is limited to the ELI5 dataset and seven language models (1B--7B parameters). The generalizability to other explanation tasks (e.g., medical, legal, or technical simplification) and larger models remains untested. Larger models (13B+) may require different architectural strategies, as they may possess sufficient parametric knowledge to reduce the benefit of explicit multi-agent decomposition.

Fourth, the multi-agent architecture introduces higher latency (29.64 seconds average) compared to single-pass prompting, which may limit applicability in real-time or interactive settings. The staged batching approach mitigates computational cost but does not reduce per-sample latency.

Finally, the adaptive synthesis strategy relies on a heuristic quality assessment rather than a learned policy. A learned strategy selector trained on retrieval quality signals could potentially improve content mixing decisions. Additionally, the creative agent uses a higher sampling temperature (0.5) than the evaluation baseline and other pipeline stages (0.1); while this is unlikely to account for LLM-judge accuracy gains, it may contribute to the observed vocabulary diversity differences and should be controlled in future work to isolate architectural from sampling effects.

\section{Conclusion}

We presented a multi-agent architecture for generating ELI5-style simplified explanations, featuring five specialized agents (breakdown, parallel reasoning and scientific, adaptive synthesis, creative) and efficient staged batching. Through comprehensive evaluation across seven language models (1B-7B parameters), 14 configurations, and approximately 66,000 samples using twelve evaluation metrics including human evaluation, we demonstrate that multi-agent orchestration consistently outperforms single-pass prompting.

GPT-4.1 evaluation shows an average 10.4\% improvement in answer accuracy, while Llama-3.3-70B evaluation shows an average 7.4\% improvement in overall quality. These gains appear across all tested models, with the largest improvements on smaller models (1B-3B) where parametric knowledge is most limited. The staged batching approach reduces inference calls from $O(N)$ to $O(1)$, while RAG integration and explicit reasoning decomposition enable superior explanation quality.

Human evaluation on 150 blind samples finds that multi-agent outputs are rated as more accurate (+10.3\% to +16.0\%) and more complete (+9.0\% to +17.8\%) than baseline outputs. However, inter-annotator agreement is moderate ($\kappa=0.162$) and evaluators diverge on simplicity, reflecting the inherently subjective nature of ELI5 quality assessment.

Importantly, we show that multi-agent explanations are 67.6\% shorter on average while achieving higher lexical overlap with references, demonstrating that the architecture produces more concise and focused explanations rather than verbose outputs. The strong negative correlation ($r = -0.99$) between word count and ROUGE-1 scores confirms that shorter, more focused explanations can achieve better alignment with reference answers.

The multi-agent architecture's advantages are most pronounced for smaller models, suggesting it is particularly valuable for resource-constrained deployments where smaller models must be used. This has important implications for edge deployment and cost-sensitive applications. Future work on hybrid architectures, specialized fine-tuning, and ablation studies may further improve performance and clarify the generalizability of these findings. Our comprehensive empirical analysis provides a foundation for informed architectural decisions in ELI5 and related explanation generation tasks.
